\subsection{Diseño del benchmark anidado}

Checkpoint~03C.2 abandona el benchmark de estimadores fijos y permite tuning disciplinado.
El track activo pasa a ser exclusivamente competition. Se conservan cuatro representaciones:
\begin{align*}
C0 &: \text{base}, && p=1400,\\
C1 &: \text{agronomic}, && p=350,\\
C2 &: \text{base+agro+empirical+discovery}, && p=2085+20,\\
C3 &: \text{base+agro+support empírico+discovery}, && p=1754+20.
\end{align*}

C3 es particularmente importante: no arrastra las 335 features empíricas, pero conserva las
primitivas configuradas necesarias para que el miner pueda explorar relaciones tipo
condición relativa/shape. Así se distingue entre usar una capa empírica completa y usar sólo
su soporte para descubrimiento.

\subsection{Nested cross-validation}

Sea $f$ un outer fold. El conjunto outer-train se divide nuevamente en tres inner folds.
Para cada configuración $\lambda\in\Lambda$:
\[
\widehat R_f(\lambda)
=
\frac{1}{3}\sum_{g=1}^{3}
\RMSE\left(
y_{V_{fg}},
\hat f_{\lambda,T_{fg}}(X_{V_{fg}})
\right),
\]
donde cualquier imputación, escalado, PCA, PLS y descubrimiento supervisado se ajusta sólo en
$T_{fg}$. Se selecciona
\[
\hat\lambda_f
=
\arg\min_{\lambda\in\Lambda}\widehat R_f(\lambda),
\]
se refita el pipeline completo sobre el outer-train y sólo entonces se predice el outer
holdout.

El inner splitter refleja el propósito del outer protocol:
\begin{itemize}
\item outer state-stratified $\rightarrow$ inner state-stratified;
\item outer municipality-grouped $\rightarrow$ inner group-by-municipality.
\end{itemize}
Esto evita seleccionar hiperparámetros bajo una geometría de validación que no corresponde a
la evaluación exterior.

\subsection{Familias de modelos}

Se evaluaron Ridge, ElasticNet, ExtraTrees, CatBoost, PLS, PCA+RBF Kernel Ridge y
PCA+polynomial-degree-2 Kernel Ridge.

\paragraph{PLS.}
PLS busca componentes latentes maximizando covarianza con $y$, por lo que puede ser útil en
high-$p$/small-$n$ con covariables correlacionadas \citep{Wold2001}. Como es supervisado,
sus componentes se ajustan exclusivamente dentro de train.

\paragraph{CatBoost.}
El boosting de árboles se ejecutó en GPU y se restringió a grids pequeños
\citep{CatBoost2018}. El objetivo era comprobar si una no linealidad potente podía extraer
señal de la capa agronómica compacta sin un barrido masivo.

\paragraph{Kernels.}
Los kernels no se materializan como datasets porque dependen de preprocessing e
hiperparámetros. Para un RBF:
\[
K_\gamma(x,z)=\exp\{-\gamma\|x-z\|_2^2\},
\]
por lo que escalar $X$ fuera del fold alteraría directamente la geometría del modelo. El mismo
principio justifica mantener PCA dentro del pipeline.

\subsection{Ranking cross-protocol}

La tabla de robustez ordena por el peor OOF RMSE entre state y grouped, y después por la media.
Los tres finalistas individuales fueron:

\begin{table}[H]
\centering
\caption{Finalistas individuales de 03C.2.}
\label{tab:cp03c2-finalists}
\begin{tabular}{llrr}
\toprule
Representación & Modelo & state OOF & grouped OOF \\
\midrule
C0 base & PLS & 0.5339 & \textbf{0.7621} \\
C3 base+agro+discovery & Ridge & 0.5331 & 0.7650 \\
C1 agronomic & CatBoost & \textbf{0.5235} & 0.7709 \\
\bottomrule
\end{tabular}
\end{table}

El orden state por sí solo sería distinto. ExtraTrees C3 obtiene 0.4972, C2 0.4996 y C1
0.5024, mejores que los tres finalistas en state, pero sus grouped RMSE son 0.8698, 0.9052 y
0.8469 respectivamente. Por tanto, el proyecto no elige simplemente el mínimo state score:
se penaliza la fragilidad espacial.

\subsection{Protocol gap}

Definimos para diagnóstico
\[
G(M)=
\RMSE_{\mathrm{grouped}}(M)
-
\RMSE_{\mathrm{state}}(M).
\]
Para C0+PLS,
\[
G\approx0.2282;
\]
para C1+CatBoost,
\[
G\approx0.2474.
\]
Incluso los mejores finalistas muestran una brecha sustancial. Esto anticipa una conclusión
fundamental de 04A: el rendimiento y los residuos tienen estructura espacial que los modelos
globales no representan completamente bajo extrapolación por municipio.

\subsection{Fracaso de PCA+RBF y PCA+Poly2}

Los kernels no lineales presentan errores muy superiores: PCA+RBF alcanza aproximadamente
3.24--4.05 RMSE, y PCA+Poly2 entre 0.72 y más de 2 dependiendo de representación/protocolo.
Esto no significa que los kernels sean universalmente inadecuados. Significa que, con sólo 138
etiquetas, el grid y preprocessing definidos no produjeron una geometría robusta. El proyecto
no siguió gastando compute en una familia que ya mostraba una brecha enorme respecto a PLS,
Ridge y árboles.

\subsection{Por qué no se hizo una búsqueda gigantesca}

Con $n=138$, aumentar el número de configuraciones evaluadas sobre los mismos folds incrementa
la posibilidad de seleccionar ruido de validación. El compute disponible no es la restricción
principal; la restricción es la cantidad de información independiente. Por ello 03C.2 utiliza
candidate grids deterministas y relativamente pequeños.

Esta política se mantiene posteriormente en 04C.1: aun disponiendo de GPU, sólo se reevalúan
tres configuraciones CatBoost heredadas cuando ya existe una hipótesis concreta de diversidad.

\subsection{Primera lectura científica}

03C.2 demuestra que:
\begin{enumerate}
\item una representación simple C0 + PLS puede ser más robusta espacialmente que ensembles
de árboles con mejor state RMSE;
\item C1 agronomic contiene suficiente estructura para que CatBoost sea competitivo;
\item la capa completa C2 no domina consistentemente a representaciones más pequeñas;
\item el modelo global todavía sufre cuando se retienen municipios enteros.
\end{enumerate}

El checkpoint constituye el núcleo de la \emph{First Modeling Delivery}, pero todavía no usa
la geometría específica de las 59 parcelas como parte explícita del proceso de selección.
